% УВОД, без номер. Ориентировъчен обем: 1 до 2 стр.
% Резюме на дейностите по курсовия проект: задача, цел, задачи, обхват, структура.
\section*{УВОД}
\label{sec:introduction}
\addcontentsline{toc}{section}{УВОД}

Обектно-ориентираният модел и програмният код, който го реализира, живеят в два
различни носителя. Моделът се описва на UML и служи за разговор между хора, а кодът се
описва на език за програмиране и служи за разговор с машина. Двата носителя се
разминават почти веднага след като работата по проекта започне, защото кодът се променя
всеки ден, а диаграмата се преначертава, когато някой се сети. Резултатът е познат:
диаграма, на която вече никой не вярва, и код, чиято структура се възстановява само с
четене.

Настоящият курсов проект разглежда това разминаване като инженерна задача с решение, а
не като организационен навик. Ако структурата на кода може да бъде извлечена машинно, а
UML моделът може да бъде записан в текстов вид, то преходът между двете посоки е
преобразувание между два формални езика и подлежи на автоматизация. Терминологичната
рамка за двете посоки е утвърдена от Chikofsky и Cross~\cite{chikofsky1990reverse}:
\term{обратно инженерство} е възстановяване на абстракция от реализация, а
\term{право инженерство} е обратният преход от абстракция към реализация. Двете заедно
образуват \term{двупосочно инженерство}, което е предметът на този проект.

\textbf{Целта} на проекта е да се проектира и реализира инструмент, наречен Blueprint,
който извлича UML модел на класове от изходен код на C++, изобразява го в текстови
нотации за диаграми, извежда от него релационна схема на данни и в обратната посока
поражда заготовки на класове и схема по подаден модел, като допълнително отчита
местата, на които моделът и кодът са се разминали.

За постигането на целта са формулирани следните \textbf{задачи}:

\begin{enumerate}[topsep=0pt,itemsep=0pt,leftmargin=*]
\item преглед на понятийния модел на UML, на градивните му елементи, на видовете
      връзки и на правилата и общите механизми, доколкото те засягат диаграмата на
      класове и нейното съответствие с изходен код;
\item анализ на съществуващите инструменти за извличане на структура от код и за
      двупосочно инженерство, и обосновка на избора на подход и на технологичен набор;
\item проектиране на вътрешния модел на системата, на границите между модулите и на
      правилата за преобразуване между клас, диаграма и релационна схема;
\item програмна реализация на извличането, на извеждането в двете текстови нотации, на
      пораждането на схема и заготовки и на проверката за разминаване;
\item изграждане на методика за експериментална проверка на верността на
      преобразуванията и на устойчивостта им при преминаване през пълен цикъл;
\item провеждане на експериментите и анализ на получените резултати.
\end{enumerate}

\textbf{Обхватът} на работата е структурната част на обектно-ориентирания модел:
класове, атрибути, операции, видимост, обхват, шаблонни класове и връзките между
класовете, заедно със съответствието им на релационна схема. Диаграмите на случаите на
употреба и на дейностите се разглеждат като контекст на модела и се използват за
описание на самата система, но не се извличат от изходен код. Извън обхвата остават
поведенческите диаграми, които изискват анализ на потока на управление, както и пълното
семантично покритие на езика C++.

\textbf{Структура на изложението.} Раздел~\ref{sec:literature} излага предметната област
и понятийния апарат на UML. Раздел~\ref{sec:alternatives} съпоставя възможните решения и
обосновава избора. Раздел~\ref{sec:design} представя проектирането на системата, а
Раздел~\ref{sec:implementation} нейната програмна реализация.
Раздел~\ref{sec:method} описва методиката за проверка, а Раздел~\ref{sec:results}
получените резултати. Заключението обобщава постигнатото и очертава по-нататъшната
работа.
